\begin{landscape}
\begin{table}[H]
\centering
\scriptsize
\caption{Perbandingan penelitian terkait ABAC, caching, dan optimasi pada IoT}
\setlength{\tabcolsep}{2pt}
\renewcommand{\arraystretch}{1.25}

\begin{tabularx}{\linewidth}{|
L{2.4cm}|
L{2.1cm}|
L{2.2cm}|
L{1.2cm}|
L{1.2cm}|
L{2.4cm}|
L{2.0cm}|
L{2.8cm}|
Y|}
\hline
\textbf{Penelitian} &
\textbf{Kelompok} &
\textbf{Domain} &
\textbf{Edge/Fog} &
\textbf{ABAC} &
\textbf{Cache/Freshness} &
\textbf{Optimasi} &
\textbf{Metrik/Evaluasi} &
\textbf{Keterbatasan} \\ \hline

\textcite{diaz_authorization_2025}
& ABAC/ authorization IoT
& IoT authorization survey
& Tidak utama
& Ya
& Terbatas
& Tidak
& Analisis komparatif model dan arsitektur otorisasi
& Membahas authorization IoT dan kebutuhan \textit{fine-grained/context-aware}, tetapi freshness atribut belum menjadi fokus utama. \\ \hline

\textcite{ullah_survey_2023}
& ABAC untuk IoT
& Blockchain-based ABAC for IoT
& Tidak utama
& Ya
& Terbatas
& Tidak
& Security/privacy, latency, overhead, scalability
& Fokus pada ABAC dan blockchain; belum membahas \textit{stale attribute} dan optimasi TTL/cache. \\ \hline

\textcite{kalaria_adaptive_2024}
& Edge/fog authorization
& Fog-based context-aware access control
& Ya
& Ya
& Tidak utama
& Tidak
& Processing time overhead, security improvement
& Fog menurunkan latency dan cloud dependency, tetapi belum fokus pada \textit{stale attribute}/freshness atribut. \\ \hline

\textcite{wang_edge-enabled_2025}
& Edge-IAM; synchronization
& Edge-enabled IAM IoT
& Ya
& Ya
& Attribute freshness dan synchronization
& Tidak
& Freshness, source trustworthiness, decision reliability, overhead
& Membahas freshness dan local authorization, tetapi belum mengoptimalkan TTL, cache size, dan cacheable attributes. \\ \hline

\textcite{penuelas-angulo_revocation_2023}
& Revocation berbasis atribut
& ABE revocation for fog-enabled IoT
& Ya
& Tidak langsung
& Revocation, bukan cache freshness
& Tidak
& User revocation, attribute revocation, cost comparison
& Membahas revocation, tetapi belum mengaitkannya dengan optimasi TTL/cache atribut untuk edge ABAC. \\ \hline

\textcite{weerasinghe_adaptive_2023}
& Context caching
& IoT context caching/CMP
& Tidak utama
& Tidak
& Context freshness
& Adaptive/RL-based caching
& Cost efficiency, QoS/QoC, freshness, latency, hit rate
& Fokus pada context/data caching, bukan cache atribut keamanan ABAC. \\ \hline

\textcite{medvedev_refresh_2023}
& IoT caching
& IoT middleware/CMP caching
& Tidak utama
& Tidak
& Freshness--latency--cost
& Cost-based caching/prefetching
& QoS, freshness, latency, operational cost
& Membahas trade-off caching, tetapi bukan otorisasi ABAC. \\ \hline

\textcite{alsadie_modified_2025}
& Metaheuristik IoT
& Fog-cloud task scheduling
& Ya
& Tidak
& Tidak
& Modified GWO
& Makespan, energy consumption
& GWO efektif untuk scheduling, tetapi belum diarahkan ke optimasi parameter ABAC. \\ \hline

\textcite{alsamarai_improved_2024}
& Metaheuristik IoT
& Fog-cloud task scheduling
& Ya
& Tidak
& Tidak
& PC-GWO
& Makespan, energy, resource utilization, cost
& Optimasi performa/resource allocation, bukan security-aware ABAC. \\ \hline

Penelitian ini
& ABAC edge security-aware optimization
& Edge IoT authorization
& Ya
& Ya
& TTL per atribut, stale hit
& Random Search, GA, GWO
& FAR, FDR, stale hit, revocation FAR, latency, PIP calls
& Simulasi terkontrol; belum diuji pada deployment fisik. \\ \hline

\end{tabularx}
\label{tab:related-work-abac-iot}
\end{table}
\end{landscape}